\documentclass{article}
\usepackage{graphicx}
\usepackage{amsmath}
\usepackage{fancyhdr}
\usepackage{parskip}
\usepackage{hyperref}

\renewcommand{\headrulewidth}{0.4pt}
\pagestyle{fancy}
\fancyhf{}
\rhead{Hireš, Jurčík, Kohoutová}
\medskip
\cfoot{\thepage}

\hypersetup{
    colorlinks=true,        
    linkcolor=blue,     
    filecolor=magenta,  
    urlcolor=cyan,
    pdfpagemode=FullScreen,
}


\begin{document}
	\begin{center}
		\section*{Is electricity production positively associated with GDP/HDI?}
	\end{center}

	\subsection*{Motivation}
		\par{}
			We examine the association between electricity production, GDP and HDI across countries.
The findings could inform policy decisions on whether to prioritise electricity infrastructure as a lever for economic development.
	\subsection*{Data description}
		\par{Electricity generation}
\\		\href{https://ourworldindata.org/grapher/electricity-generation.csv?v=1&csvType=full&useColumnShortNames=true}{Original dataset}
\\4 columns : Entity (Country name), Code (Country code), Year, Total Electricity (in TWH)
\\7914 rows including different countries and world regions across different years 
\\Date range: 1985–2025
\\Unit: Terawatt hours (TWH)
		\par{GDP}
		\\\href{https://ourworldindata.org/grapher/gdp-penn-world-table.csv?v=1&csvType=full&useColumnShortNames=true}{Original dataset}
\\4 columns: Entity, Code, Year, RGDP (Real gross domestic product in USD)
\\10908 rows including different countries across different years 
\\This data is expressed in international-\$ at 2021 prices, using a multiple benchmark approach that incorporates PPP estimates from all available benchmark years
\\Date range: 1950–2023
		\par{HDI}
		\\\href{https://ourworldindata.org/grapher/human-development-index.csv?v=1&csvType=full&useColumnShortNames=true}{Original dataset} 
\\ 5 columns: Entity, Code, Year, hdi\_\_sex\_total (Human development), \_region (world region)
\\6605 rows including different countries and world regions, income classes of countries across different years 
\\The Human Development Index (HDI) is the geometric mean of the three dimensions: health, measured life expectancy at birth; education, measured by the expected and mean years of schooling; and material wellbeing, measured by GNI per capita (PPP\$). Each indicator is normalized to a scale from 0 to 1 using fixed minimum and maximum values. 
\\Date range: 1990–2023
	\subsection*{Data Cleaning}
		\par{}
		We removed world regions from the data by keeping only entries with country code.
We also removed owid\_regions, because they do not affect our thesis. 
We excluded Lesotho as it has recorded 0 TWh in the source data, which prevents log transformation.
We filtered countries with data entries in 2000 – 2023, to ensure consistent comparison between years. This specific range is because HDI and GDP entries are only until year 2023 and many developing countries do not have any entries before the year 2000.
		\par{} 
		This step removed following countries from our analysis:
Antigua and Barbuda, Belarus (former USSR state reporting late), Bermuda, Cayman Islands, Montenegro, South Sudan (recent independence), Ukraine (former USSR state reporting late), Venezuela (Data suppression).
		\par{}
		So in the end we were left with two important datasets final\_start\_HDI and final\_end\_HDI, both containing 172 countries with 9 columns (code ,year, total\_generation\_\_twh, hdi\_\_sex\_total, rgdpo, change\_gdp, elec\_intensity, predicted\_log\_twh, residual)
		\subsection*{Descriptive statistics}
		\par{}
		Electricity production and real GDP both exhibit strong right-skew: the mean exceeds the median by a factor of 11 for electricity and 7.6 for GDP. The standard deviation in each case dwarfs the mean. A handful of large producers (China, USA, India) pull the average upward while the typical country sits near the median. This dispersion across orders of magnitude is what motivates the log transformation in a later section.
		\par{}
		HDI shows a different pattern. The mean and median are nearly equal and the distribution is mildly left-skewed - more countries cluster at the upper end of the scale, with a thin tail toward lower development. The bounded nature of HDI limits the scope for variation among developed countries and weakens any linear relationship with electricity at the upper range.
		\par{}
		HDI is constructed as a per-capita-style composite, so unlike GDP and electricity totals it does not scale with population - which explains why its distribution is more symmetric.
		\par{}
		Electricity intensity (TWh per USD of real GDP) ranges across two orders of magnitude.  The standard deviation is roughly equal to the mean, indicating that countries differ substantially in how much electricity their economies produce per unit of output - reflecting differences in industrial composition, climate, energy efficiency, and the role of electricity-intensive sectors such as aluminum smelting or hydroelectric exports.
	\newpage
	\subsection*{Visualisations}
		\par{}
		total\_graph\_rgdpo
		\\\begin{center}
			\\\includegraphics[scale = 0.5]{total_graph_rgdpo.pdf}
		\end{center}
		\par{}
		We initially plotted real GDP against electricity production on a linear scale.
		\par{}
		There is a positive correlation between electricity production and real GDP with $R^2 = 0.906$. However it is not well interpretable as the whole graph is scaled by the biggest observations, China and the United States, whose values lie an order of magnitude above all others. The remaining majority of countries collapsed into a dense cluster near the origin.		
		\par{}
		hist\_electricity
		\\\begin{center}
			\\\includegraphics[scale = 0.5]{hist_electricity.pdf}
		\end{center}
		\par{}
		On the $log10$ axis, the distribution of electricity production is approximately symmetric and centered on the median, indicating that the underlying values are roughly log-normal. The mean sits to the right of the median, which reflects right skewness in the original linear scale. A non-log histogram would compress almost all 172 countries into a single bin near zero, with a long tail extending to a handful of large producers; this is the reason the log transform is used. The symmetry on the log scale is not coincidental but the expected pattern for quantities that span several orders of magnitude and arise from multiplicative growth processes. Large producers such as China and the United States raise the mean substantially while leaving the median almost unchanged, because the median is fixed by the rank position of the middle country and is insensitive to extreme values in the tail.
		\par{}
		start\_graph end\_graph
		\\\begin{center}
			\\\includegraphics[scale = 0.5]{start_end_graph.pdf}
		\end{center}
				\par{}
		We used a log-log graph to see the relationship better for all countries. The correlations are very strong as well. If we compare the year 2000 and 2023, we can see that the correlation increased as the countries follow the trend more tightly. We can also see general economic progress, as countries moving upwards and to the right.
		\newpage
		\par{}
		graph\_hdi
		\\\begin{center}
			\\\includegraphics[scale = 0.5]{total_graph_hdi.pdf}
		\end{center}
		\par{}
		Log-transformed correlation between electricity production and HDI is $r = 0.568$ ($r^2 = 0.322$), indicating a moderate positive relationship. From this graph we can see that the correlation between HDI and electricity production is 
smaller and the linear regression line is flatter. HDI is bounded on [0,1] by construction, so countries near the upper limit have little room to move regardless of how much their electricity output grows; this flattens the slope at the top of the distribution. A secondary economic explanation is that once a country reaches a certain industrial level, additional electricity contributes less to HDI gains. Beyond that point, HDI depends more on social structures than on industrial output.
		\par{}
		Aruba, British Virgin Islands, Macao, Montserrat, and Taiwan lack 2023 HDI data and are excluded from analysis.


	\subsection*{Central findings}
		\par{}
		We initially plotted real GDP against electricity production on a linear scale. There is a positive correlation with $R^2 = 0.906$, but this graph is dominated by two observations, China and the United States, whose values lie an order of magnitude above all others (these two points effectively determine the slope of the regression). The remaining countries collapse into a dense cluster near the origin and cannot be distinguished on this scale.
		\par{}
		The distribution of electricity production is highly right-skewed on the linear scale. Most countries cluster at low to moderate output, while a small number produce volumes orders of magnitude greater. This Pareto-like shape reflects the concentration of industrial capacity in a small group of large economies. The log transformation addresses this by placing countries on a scale where proportional rather than absolute differences are visible, which is the appropriate framing when values span several orders of magnitude.
		\par{}
		While the overall relationship is strong, individual countries deviate systematically from the trend. We computed residuals from the log-log regression to identify countries whose electricity production differs substantially from what their GDP level would predict.
		\par{}
		Below the trend are two groups: low-income countries with weak grid infrastructure (Chad, Sierra Leone, Niger, Guinea-Bissau, Djibouti), where GDP exists despite limited electrification; and small economies that import electricity from neighbouring countries (Luxembourg, Macao, Palestine).
		\par{}
		Above the trend are also two groups: hydroelectric exporters (Bhutan, Iceland, Laos, Paraguay, Tajikistan); and resource-intensive small economies (Bahrain, Kuwait).
		\par{}
		These outliers reveal that the GDP-electricity relationship is affected by industrial composition, resource endowment, and trade in electricity.
		\begin{table}[!htbp] \centering
\caption{Top 5 negative and top 5 positive residuals from the 2023 log-log regression of GDP on electricity production. Residual = log10(actual electricity) - log10(predicted electricity). Factor = $10^{\text{residual}}$: a factor of 0.09 means the country produces 9\% of its predicted electricity, a factor of 10.96 means nearly 11x predicted.}
\label{tab:residuals}
\begin{tabular}{@{\extracolsep{5pt}} cccccc}
\\[-1.8ex]\hline
\hline \\[-1.8ex]
Direction & Country & log10(elec) & log10(GDP) & Residual & Factor \\
\hline \\[-1.8ex]
Below & Chad & $-0.44$ & $10.62$ & $-1.06$ & $0.09$ \\
Below & Sierra Leone & $-0.68$ & $10.38$ & $-1.04$ & $0.09$ \\
Below & Macao & $-0.18$ & $10.66$ & $-0.84$ & $0.14$ \\
Below & Somalia & $-0.38$ & $10.41$ & $-0.76$ & $0.17$ \\
Below & Guinea-Bissau & $-1.10$ & $9.72$ & $-0.74$ & $0.18$ \\
Above & Barbados & $0.04$ & $9.49$ & $0.65$ & $4.47$ \\
Above & Paraguay & $1.72$ & $11.01$ & $0.69$ & $4.90$ \\
Above & Laos & $1.72$ & $10.81$ & $0.90$ & $7.94$ \\
Above & Iceland & $1.30$ & $10.39$ & $0.93$ & $8.51$ \\
Above & Bhutan & $1.05$ & $10.05$ & $1.04$ & $10.96$ \\
\hline \\[-1.8ex]
\end{tabular}
\end{table}
	\subsection*{Limitations}
		\par{}
		We found a positive correlation between electricity production and real GDP.
The causation remains unanswered, we suspect that causation runs in both directions: GDP funds grid infrastructure and electricity enables economic activity.
		\par{}
		However, we did not account for the population of the country at all. It might very likely be that bigger population causes both increase in GDP and in electricity production. In order to understand the deeper relationship between these two, we would need to compute everything per capita, which could be an interesting topic for further study.

	\subsection*{Collaboration}
		\par{}
		We brainstormed the idea together, afterwards Max and Kryštof did the Data Cleaning, Descriptive statistics, Visualizations. Kristýna then wrote the Central findings, Limitations and Slides.
\end{document}
